\documentclass[11pt]{article}
\usepackage{fontspec}
\usepackage[french,english]{babel}
\usepackage{amsmath}
\usepackage{amssymb}
\usepackage{amsthm}
\usepackage{mathrsfs}
\usepackage{tikz-cd}
\usepackage{enumitem}

\newcommand{\printedpage}[1]{\marginpar{\footnotesize p.~#1}}
\newcommand{\msp}[1]{\marginpar{\footnotesize [#1]}}
% circled labels of the conditions, as printed
\newcommand{\circled}[1]{\tikz[baseline=(c.base)]\node[draw,circle,inner sep=1pt](c){#1};}
% "defined as" sign, as printed
\newcommand{\defeq}{\overset{\text{def}}{=}}
% tilde printed after a closing parenthesis
\newcommand{\postilde}{\tilde{}}
\newcommand{\Fl}{\mathrm{Fl}}
\newcommand{\Cof}{\mathrm{Cof}}

\begin{document}

\printedpage{34}
% head: A. Grothendieck : Dérivateurs, Ch. XVII, Édition des Universités de Montpellier II et Paris VII
% foot: Version 9 février 2019
The \textit{strict} hypothesis for a category with cofibrations seems fairly innocuous to me, in most cases known to me, it is satisfied. But even with this hypothesis, I see no way to prove the existence of $u_!$ for \textit{every} arrow of Cat, and this is of course a very awkward situation. One would very much like to isolate conditions that ensure the validity of Der 4 bis and of Der 5 bis (existence of $u_!$ for
\msp{page 71}
\textit{every} arrow of Cat, and way of computing the fibers \dots).

\section*{4\quad Special categories with cofibrations}

\textbf{4.1.} I now suppose that
\begin{itemize}[leftmargin=3em]
\item[\circled{a}] $\mathcal{M}$ is stable under filtering $\varinjlim$, and
\item[\circled{b}] that its objects are accessible,
\end{itemize}
so that we are under the conditions for applying the factorization theorem (XV, \S1). We suppose:
\begin{itemize}[leftmargin=3em]
\item[\circled{c}] given $W, C \subset \Fl(\mathcal{M})$, satisfying the conditions $\Cof_1$, $\Cof_2$ of page 39. For $\Cof_3$, cf. below.
\item[\circled{d}] $C$, $TC \defeq C \cap W$ \textit{stable under transfinite composition} (this is the case if the structure comes from a closed Quillen structure $(W, C, F)$).
\end{itemize}
\[
TF \defeq C_* \;\subset\; \Fl(\mathcal{M})\ , \tag*{(4.1)}
\]
it follows from the factorization theorem that every arrow $f$ of $\mathcal{M}$ factors as
\[
f \;=\; pi\ , \qquad p \in TF,\ i \in C\ ,
\]
\textit{provided} we make the hypothesis
\begin{itemize}[leftmargin=3em]
\item[\circled{e}] \mbox{}
\[
\exists\, C_0 \subset C,\ \textit{small} \text{ subset of } C, \text{ such that } C \subset \tilde{C}_0, \text{ i.e. } TF = (C_0)_*. \tag*{(4.3)}
\]
\end{itemize}
Under these conditions, if we suppose
\begin{itemize}[leftmargin=3em]
\item[\circled{f}] \mbox{}
\[
TF \subset W\ , \tag*{(4.4)}
\]
\end{itemize}
\msp{page 72}
then the factorization axiom $\Cof_3$ is satisfied. Conversely, if $\Cof_3$ is satisfied, and if $W$ is stable under direct factors, then we have 4.4. Thus if $W$ is stable under direct factors, $\Cof_3$ is equivalent to (4.4).

\textbf{4.2.} It seems to me that these conditions, in practice, will hardly ever be fulfilled, unless $W$ (which is given in the very first place, satisfying $\Cof_0$ \textsuperscript{(}\footnote[18]{i.e. weak saturation.}\textsuperscript{)}), is not stable under $\varinjlim$ that are
\printedpage{35}
% head: A. Grothendieck : Dérivateurs, Ch. XVII, Édition des Universités de Montpellier II et Paris VII
% foot: Version 9 février 2019
filtering. At that point, to construct the cofibration structure, we start from $C_0$ a small subset of $\Fl(\mathcal{M})$, taken so that we have
\[
(C_0)_* \;\subset\; W\ , \tag*{(4.5)}
\]
we take for $C$ at choice the saturation $C'_0$ of $C_0$ under cobase change and transfinite composition, or $C = \tilde{C}_0 =$ set of direct factors of $C'_0$. The conditions $\Cof_1$ (stabilities of $C$), $\Cof_3$ (factorization) are satisfied automatically, and it remains to take care of $\Cof_2$, which comes down to the stability of $C{\cap}W$ under cobase change (stability under composition and that it contain the isomorphisms, being immediate). This is a delicate condition, and the only standard way I see to ensure it, is by supposing
\[
C_0 \;\subset\; \Cof_W\ ; \tag*{(4.6)}
\]
\msp{page 73}
the stability of $W$ under filtering $\varinjlim$ (hence also under direct factors) then ensures that we also have
\[
C \;\subset\; \Cof_W\ , \tag*{(4.7)}
\]
and consequently
\[
\left\{
\begin{aligned}
& C \cap W \subset \Cof_W \cap W = W^{\mathrm{univ}}, \\
& \text{hence }\ TC \defeq C \cap W = C \cap W^{\mathrm{univ}}\ ,
\end{aligned}
\right. \tag*{(4.8)}
\]
which ensures the stability of $C \cap W$ under cobase change. The stability of $TC$ under transfinite composition follows from that of $C$, and from the stability of $W$ under transfinite composition.

\textbf{4.3.} I return to the general conditions\circled{a} to\circled{f} of 4.1, without necessarily supposing $W$ stable under filtering $\varinjlim$. I now want to extend the cofibration structure to the $\mathcal{M}(I)$, $I$ arbitrary in Cat. For this we do not (I think) have to use the developments of sections 2, 3, but we will make the connection with these later on.

Let us denote by $I_0$ the discrete category having the same objects as $I$, whence a canonical functor
\[
v = v_I \ :\ I_0 \longrightarrow I\ . \tag*{(4.9)}
\]
The idea is to define the structure of $\mathcal{M}(I)$ via the evident structure of $\mathcal{M}(I_0) \simeq \mathcal{M}^{I_0}$, by playing on the pair of adjoint functors
\[
\begin{tikzcd}
\mathcal{M}(I_0) \arrow[r, shift left, "\varphi = v_!"] & \mathcal{M}(I) \arrow[l, shift left, "\psi = v^*"]
\end{tikzcd}
\ .
\]
\msp{page 74}
We are going to present things in a slightly more general way.

\textbf{Definition 4.4.} I call a \textit{special} cofibration structure a structure $(\mathcal{M}, W, C)$ satisfying conditions a) to f) of 4.1, which I reformulate:
\begin{itemize}[leftmargin=3em]
\item[\circled{a}] $\mathcal{M}$ stable under small direct limits, and the objects of $\mathcal{M}$ are accessible.
\printedpage{36}
% head: A. Grothendieck : Dérivateurs, Ch. XVII, Édition des Universités de Montpellier II et Paris VII
% foot: Version 9 février 2019
\item[\circled{b}] $W$ is weakly saturated (i.e. $\Cof_0$ of p. 39), $C$ stable under cobase change and under composition, $TC \defeq C{\cap}W$ stable under cobase change (and under composition automatically).
\item[\circled{b$'$}] $C$ and $TC$ are stable even under composition that is \textit{transfinite.}
\item[\circled{c}] $\exists$ small subset $C_0 \subset C$ such that $C \subset \tilde{C}_0$ ($\Longrightarrow \tilde{C} = \tilde{C}_0$).
\item[\circled{d}] $TF \defeq C_* \subset W$ (which implies, together with\circled{a} to\circled{c}, the factorization axiom $\Cof_3$ of page 39).
\end{itemize}

\textbf{4.5.} We are going to give a general recipe for extending such a structure, when we have a pair of adjoint functors
\[
\begin{tikzcd}
\mathcal{M} \arrow[r, shift left, "v_!"] & \mathcal{M}' \arrow[l, shift left, "v^*"]
\end{tikzcd}
\ ,
\]
$\mathcal{M}$ being equipped with a special cofibration structure $(W, C)$, and $\mathcal{M}'$ satisfying the preliminary condition\circled{a}. We will then define the subsets $W'$, $C'$ of $\Fl(\mathcal{M}')$
\msp{page 75}
by
\[
\left\{
\begin{aligned}
W' \;&=\; (v^*)^{-1}(W) \\
C' \;&=\; \text{closure stable under cobase change} \\
&\phantom{\;=\;}\ \ \text{and transfinite composition of } v_!(C),
\end{aligned}
\right. \tag*{(4.5.1)}
\]
thus
\[
v_!(C) \;\subset\; C' \;\subset\; (v_!(C))\postilde \;=\; \tilde{C}'\ . \tag*{(4.5.2)}
\]
I say that $(W', C')$, and at choice $(W', \tilde{C}')$, is a special cofibration structure on $\mathcal{M}'$. We are going to check the various conditions of 4.4. Condition a) is already assured.

b) \textsuperscript{(}\footnote[19]{\textbf{N.B.} It is better to start by proving c).}\textsuperscript{)}. The weak saturation of $W'$ follows from that of $W$, via (4.5.1). The stability of $C'$ under cobase change and under composition (even transfinite) is assured by definition of $C'$, and likewise for $\tilde{C}'$ (being Q-closed on the left). Stability of $TC' \defeq C'{\cap}W'$ under cobase change? It suffices to see that if we have a cocartesian square
\[
\begin{tikzcd}
X'_0 \arrow[r, "i'"] \arrow[d] \arrow[dr, phantom, "\text{\scriptsize cocart.}"] & X'_1 \arrow[d] \\
X'_2 \arrow[r, "j'"] & X'_3
\end{tikzcd}
\]
in $\mathcal{M}'$, with $i'$ in $TC' = C' \cap W'$, then $j' \in W'$ (for we already know $j' \in C'$), i.e.
\marginpar{$||$}% marginal sign || printed in the left margin of this line
$v^*(j') \in W'$. We are going to suppose that $v^*$ \textit{commutes with amalgamated sums,} and we are thus reduced to proving that $v^*(i') \in TC$. Thus % [sic] $v^*(j') \in W'$ -> $v^*(j') \in W$, since $W' = (v^*)^{-1}(W)$
\printedpage{37}
% head: A. Grothendieck : Dérivateurs, Ch. XVII, Édition des Universités de Montpellier II et Paris VII
% foot: Version 9 février 2019
\msp{page 76}
we must prove the

\textbf{Lemma 4.5.2.} \textit{Under the conditions} (4.5.3) \textit{and} (4.5.4) \textit{below, we have}
\[
v^*(C') \;\subset\; C\ , \qquad v^*(TC') \;\subset\; TC\ ,
\]
\textit{and if $C = \tilde{C}$, we have $v^*(\tilde{C}') \subset C$, $v^*(T\tilde{C}') \subset TC$.}

For $v^*(C') \subset C$, we must first prove
\[
v^*(v_!(C)) \;\subset\; C\ , \tag*{(4.5.3)}
\]
and I see no way -- we have to take it as a hypothesis. Moreover, that $v^*$ commutes with arbitrary $\varinjlim$ (cobase change already admitted -- now we need filtering direct limits)
\[
v^* \text{ commutes with small } \varinjlim \text{ of any kind} \tag*{(4.5.4)}
\]
to conclude from (4.5.3) that $v^*(C' =$ closure of $v_!(C)$ under cobase change and transfinite composition) is also in $C$. Passing to $v^*(TC') \subset TC$, by what precedes, $v^*(C') \subset C$, it suffices to prove $v^*(TC') \subset W$, but $v^*(TC') \subset v^*(W') \subset W$ by definition of $W'$.

Let us suppose $C = \tilde{C}$, i.e. $C$ stable under direct factors, and let us prove $v^*(\tilde{C}') \subset C$, $v^*(T\tilde{C}') \subset TC$. By c) below, $\tilde{C}'$ is formed of the direct factors
\msp{page 77}
of elements of $C'$, hence $v^*(C') \subset C$ implies $v^*(\tilde{C}') \subset C$ since $C$ stable under direct factors.

Let us finally prove $v^*(T\tilde{C}') \subset TC \defeq C \cap W$. As $T\tilde{C}' = \tilde{C}' \cap W'$, this follows from $v^*(\tilde{C}') \subset C$, and from $v^*(W') \subset W$.

Let us prove b$'$). The stability of $C'$ and of $\tilde{C}'$ under transfinite composition is assured, it remains to prove the stability of $TC' = C' \cap W'$ and of $T\tilde{C}' = \tilde{C}' \cap W'$. It suffices to see that the transfinite composite $f$ in question is in $W'$, hence that $v^*(f) \in W$, but this follows from the fact that $v^*$ commutes with transfinite composition by (4.5.4), and that $v^*$ takes $TC'$ and $T\tilde{C}'$ into $TC$.

\textit{Let us prove} c) for $C'$ (hence also for $\tilde{C}'$). Let $C_0$ be a small subset of $C$ such that $C \subset \tilde{C}_0$, i.e. $C$ is contained in the closure of $C_0$ under the three operations of cobase change, of transfinite composition, of direct factors. We then see that $v_!(C)$ is in the closure of $v_!(C_0)$ under the operations, hence also $C'$, and we can then take $C'_0 = v_!(C_0)$.

\textit{Let us finally prove} d), it is the same assertion for $C'$ and for $\tilde{C}'$, since
\[
C'_* \;=\; \tilde{C}'_*\ (\ \defeq\ TF'\ )\ ,
\]
\msp{page 78}
we must therefore prove
\[
TF' \defeq C'_* \;\subset\; W'\ .
\]
This will follow from the
\printedpage{38}
% head: A. Grothendieck : Dérivateurs, Ch. XVII, Édition des Universités de Montpellier II et Paris VII
% foot: Version 9 février 2019

\textbf{Lemma 4.5.3.} \textit{We have $TF' = (v^*)^{-1}(TF)$.}

This is formal from the definition of $TF'$ as
\[
TF' \;=\; (v_!(C))_*\ ,
\]
of $TF$ as
\[
TF \;=\; C_*\ ,
\]
and from the adjunction of $v_!$, $v^*$. As $TF \subset W$, we conclude
\[
TF' \;=\; (v^*)^{-1}(TF) \;\subset\; (v^*)^{-1}(W) = W'\ ,
\]
q.e.d.

Thus we have proved this:

\textbf{Proposition 4.6.} \textit{Let us consider a pair of adjoint functors}
\[
\begin{tikzcd}
\mathcal{M} \arrow[r, shift left, "v_!"] & \mathcal{M}' \arrow[l, shift left, "v^*"]
\end{tikzcd}
\ ,
\]
\textit{where $\mathcal{M}$ is equipped with a cofibration structure that is} special $(W, C)$\textit{. We further suppose the following conditions satisfied.}
\[
\begin{aligned}
&1^\circ)\quad \mathcal{M}' \textit{ stable under small } \varinjlim\ \textit{, its elements are accessible.} \\
&2^\circ)\quad v^* \textit{ commutes with small } \varinjlim\ . \\
&3^\circ)\quad v^*(v_!(C)) \subset C.
\end{aligned} \tag*{(4.6.1)}
\]
\msp{page 79}
\textit{Let us consider the subsets $W'$, $C'$ of} $\Fl(\mathcal{M}')$ \textit{defined by}
\[
\left\{
\begin{aligned}
W' \;&=\; (v^*)^{-1}(W), \\
C' \;&=\; \textit{closure of } v_!(C) \textit{ under cobase change and under composition} \\
&\phantom{\;=\;}\ \textit{that is transfinite.}
\end{aligned}
\right. \tag*{(4.6.2)}
\]
\textit{Then $(\mathcal{M}', W', C')$ is a special cofibration structure. If moreover $C = \tilde{C}$, i.e. $C$ is stable under direct factors, then $(\mathcal{M}', W', \tilde{C}')$ is also a special cofibration structure} \textsuperscript{(}\footnote[20]{\textbf{N.B.} This is a general fact for special categories with cofibrations, and should be made explicit below (cf. 5.1). The hypothesis on $C$ is unnecessary here (or rather, it suffices instead that $W$ be stable under direct factors \textit{and} under transfinite composition, which is not a more cheering hypothesis \dots).}\textsuperscript{)}.

\textbf{4.7.} I would like, with a view to statements of type 3.7 (establishing the existence of functors of type $u_!^W$), to have a hold on $TC' = C'{\cap}W'$, and on $T\tilde{C}' = \tilde{C}'{\cap}W'$, to be precise, I would like to have the relation
\[
v_!(TC)\ \subset TC' \subset v_!(TC)\postilde\ , \tag*{(4.7.1)}
\]
or even the slightly stronger relation
\[
\underbrace{v_!(TC)\ \subset TC'}_{} \ \overset{\text{tautol.}}{\subset}\ \underbrace{T\tilde{C}' \;\subset\; (v_!(TC))\postilde}_{}\ , \tag*{(4.7.2)}
\]

\end{document}
